\documentclass[10pt,a4paper]{amsart}
\usepackage[margin=19mm]{geometry}
\usepackage{amsmath,amssymb,mathtools}
\usepackage[scr=rsfs,cal=euler]{mathalfa}
\usepackage{xfrac}
\usepackage[unicode,hidelinks,bookmarks=false]{hyperref}
\newcommand{\ie}{i.e.\ }
\newcommand{\cf}{cf.\ }
\newcommand{\resp}{respectively\ }
\newcommand{\Subsection}{\subsection}
\providecommand{\nobreakdash}{\nobreak\hbox{-}}
\setlength{\emergencystretch}{2em}
\multlinegap0pt
\usepackage{tikz}
\usetikzlibrary{cd}
\usepackage{amssymb}
\usepackage{csquotes}
\usepackage{braket}
\usepackage{xcolor}
\usepackage{tikz}
\usepackage{float}
\usepackage[shortlabels]{enumitem}
\newtheorem{theorem}{Theorem}
\newcommand{\afa}{\mathsf{AFA}}
\DeclareMathOperator{\con}{Con}
\DeclareMathOperator{\uni}{\mathbf{V}}
\newcommand{\zfc}{\mathsf{ZFC}}
\title{A new kind of anti-foundation axioms}
\author{Daheng Ju}
\date{Source: arXiv:2608.21682v1 (CC BY 4.0)}
\begin{document}
\maketitle

\noindent\emph{Source and licence.} A new kind of anti-foundation axioms. Version arXiv:2608.21682v1 (CC BY 4.0), distributed by arXiv under the Creative Commons Attribution 4.0 International licence, http://creativecommons.org/licenses/by/4.0/, as recorded in the arXiv licence metadata for this exact version and shown on the abstract page https://arxiv.org/abs/2608.21682v1. Full licence notice: \texttt{licenses/arxiv-2608.21682-CC-BY-4.0.txt}.\par\smallskip

\noindent\textit{Editorial scope.} Counted unit: the article's theorem theorem (source0.tex lines 351--353). The article's complete original proof of it is delivered with it (source0.tex lines 355--397, one \texttt{proof} environment of 43 lines). No mathematics is written, repaired or re-derived: the statement and the proof are the article's own lines, in the article's own notation, and no retained line is substituted or re-flowed.  No further source-local context is needed: every cross-reference of every retained line resolves inside the two delivered spans.  Nothing was omitted from any retained span, so the package carries no omission record.  No retained line contains a citation, so the wrapper carries no bibliography and no bibliography entry.  The retained lines contain the article's own diagram code, which the wrapper typesets with the standard tikz packages; no image file is delivered and the package declares no asset dependency.  Editorial formatting only, in addition: an AMS \texttt{amsart} preamble that restates the article's own declarations and macros used by the retained lines, namely the environments \texttt{theorem} on separate counters, together with the standard packages those lines need.  The retained environments print in this document as Theorem 1 in order of appearance, whereas the article prints the same items with the numbering of its own full document.  The complete native source of the article as distributed in the arXiv e-print for this exact version is bundled unchanged as \texttt{sources/208268/source0.tex}.
\begin{theorem}
    (finitary arithmetic) There exists a graph relation $\sim$ such that $\zfc^-\vdash$``$\sim$ is regular'', $\zfc^-+\afa^\sim$ is consistent relative to $\zfc^-$, and for every $\sim'$ such that $\zfc^-\vdash$``$\sim'$ is a regular bisimulation'', (if $\zfc^-$ is consistent, then) $\zfc^-\not\vdash\afa^\sim\Leftrightarrow\afa^{\sim'}$.
\end{theorem}

\begin{proof}
    Denote by $Q_n$ the graph consisting of $n$ ($\geq 1$)-many vertices with self-loops $a_1\to\dots\to a_n$ having $a_1$ as point.
    
\begin{figure}[htbp]
\[\begin{tikzcd}
	{a_1} \\
	\\
	{a_2} \\
	\\
	{a_3}
	\arrow[from=1-1, to=1-1, loop, in=325, out=35, distance=10mm]
	\arrow[from=1-1, to=3-1]
	\arrow[from=3-1, to=3-1, loop, in=325, out=35, distance=10mm]
	\arrow[from=3-1, to=5-1]
	\arrow[from=5-1, to=5-1, loop, in=325, out=35, distance=10mm]
\end{tikzcd}\]
\caption{$Q_3$}
\end{figure}

    For $n\geq2$, define $\cong_{Q_n}^*$ as follows: $G\cong_{Q_n}^*G'$ iff $G\cong^*G'$ or $G,G'\in\{Q_1,Q_1^*,Q_n,Q_n^*\}$ up to isomorphism.
    
    Obviously, $\cong_{Q_2}^*$ is a regular bisimulation. For $n\geq3$, $\cong_{Q_n}^*$ is regular, and below we prove that $\cong_{Q_n}^*$ satisfies the other requirements.

    First, we prove that $\zfc^-+\afa^{\cong_{Q_n}^*}$ is consistent relative to $\zfc^-$.
    
    By Corollary 3, it's sufficient to prove under $\zfc^-$ that $\cong_{Q_n}^*$ has the set property and is absolute for $\uni_c^{\cong_{Q_n}^*}$.
    The latter is obvious. Below we prove the former, and it's sufficient to prove that $\cong_{Q_n}^*$ satisfies (III) in Theorem 3.
    
    (i): If $G\not\cong G'$ are both $\cong_{Q_n}^*$-extensional, then they are both $\cong^*$-extensional, so $G\not\cong^*G'$.
    Assume the contrary that $G\cong_{Q_n}^* G'$, comparing $\cong^*$ and $\cong_{Q_n}^*$, we have that $G,G'\in\{Q_1,Q_1^*,Q_n,Q_n^*\}$ up to isomorphism. Since $G,G'$ are both $\cong_{Q_n}^*$-extensional, and among $Q_1,Q_1^*,Q_n,Q_n^*$, only $Q_1$ is $\cong_{Q_n}^*$-extensional, we have that $G,G'\cong Q_1$, contrary to $G\not\cong G'$!
    Thus, $G\not\cong_{Q_n}^* G'$.
    
    (ii): If $G$ has a point having no parents and is quasi-$\cong_{Q_n}^*$-extensional, then it is quasi-$\cong^*$-extensional, so for every $a\in G\backslash\{p_G\}$, $G\cong^* G[a]\Rightarrow G\approx G[a]$.
    Assume the contrary that there exists some $a\in G\backslash\{p_G\}$ such that $G\cong_{Q_n}^* G[a]$ and $G\not\approx G[a]$, comparing $\cong^*$ and $\cong_{Q_n}^*$, we have that $G,G[a]\in\{Q_1,Q_1^*,Q_n,Q_n^*\}$ up to isomorphism. Since $G$ has a point having no parents and is quasi-$\cong_{Q_n}^*$-extensional, we have that $G\cong Q_1^*$, so $G[a]\cong Q_1$, contrary to $G\not\approx G[a]$!
    Thus, for every $a\in G\backslash\{p_G\}$, $G\cong_{Q_n}^* G[a]\Rightarrow G\approx G[a]$.
        
    Second, we prove that for every $\sim$ such that $\zfc^-\vdash$``$\sim$ is a regular bisimulation'', (if $\zfc^-$ is consistent, then) $\zfc^-\not\vdash\afa^{\cong_{Q_n}^*}\Leftrightarrow\afa^{\sim}$.
    
    Working in $\zfc^-$, we prove that $\afa^{\cong_{Q_n}^*}\Rightarrow\neg\afa^\sim$. Note that $Q_2$ is $\cong_{Q_n}^*$-extensional and $Q_n$ is not $\cong_{Q_n}^*$-extensional.
    Assume the contrary that $\afa^{\cong_{Q_n}^*}\land\afa^\sim$, then $Q_2$ is $\sim$-extensional and $Q_n$ is not $\sim$-extensional, i.e., there exist some $1\leq i<j\leq n$ such that $Q_i\sim Q_j$. By repeatedly applying the condition that $\sim$ is a regular bisimulation, we have that $Q_1\sim Q_2$, so $Q_2$ is not $\sim$-extensional, a contradiction! Thus, $\afa^{\cong_{Q_n}^*}\Rightarrow\neg\afa^\sim$.
    
    Back to finitary arithmetic, we have that $\zfc^-+\afa^{\cong_{Q_n}^*}\vdash\neg\afa^\sim$, so $\con(\zfc^-)$ implies $\con(\zfc^-+\afa^{\cong_{Q_n}^*})$, which in turn implies $\con(\zfc^-+\afa^{\cong_{Q_n}^*}\not\Leftrightarrow\afa^{\sim})$.
\end{proof}

\end{document}
